\section{Generalized Killing--Eisenhart and Levi-Civita conditions}

The proofs of the main results of our paper, that is Theorems \ref{maintheorem1}, \ref{maintheorem2} and \ref{maintheorem3}, make use of generalizations to Painlev\'e metrics of the classical Killing--Eisenhart equations and Levi-Civita separability conditions which hold for St\"ackel metrics (see for example \cite{BCR1-2002, BCR2-2002, Eis1934, KaMi1980, KaMi1981, Per1990}). We present these generalizations in the form of the following two lemmas, beginning with the Killing--Eisenhart equations. Thus in analogy with the St\"ackel case, we introduce the quantities
\begin{gather}\label{defrho}
\rho_{\beta \gamma}:=\frac{s^{\gamma \beta}}{s^{\gamma 1}}.
\end{gather}
Note that by the assumption (\ref{StackelPositiveness}), we have $s^{\gamma 1}\neq 0$.
The following lemma gives the generalization to the case of Painlev\'e metrics of the Killing--Eisenhart equations given in \cite{BCR1-2002,BCR2-2002,Eis1934,KaMi1980,KaMi1981} for St\"ackel metrics:
\begin{Lemma} \label{KE}
We have, for all $1\leq \beta, \delta, \gamma \leq r$, the identities
\begin{gather}\label{genKE}
\partial_{j_{\gamma}}(\rho_{ \beta \delta})= (\rho_{\beta \gamma}-\rho_{\beta \delta })\left(\partial_{j_{\gamma}}{\log \frac{s^{\delta 1}}{\det S}}\right).
\end{gather}
\end{Lemma}
We will refer to \eqref{genKE} as the {\emph{generalized Killing--Eisenhart equations.}}
\begin{proof}The Poisson bracket relations (\ref{PoissonCommute}) imply
\begin{gather}\label{Painlevecomm}
\sum_{\pgamma=1}^{\lgamma}\big(\big(\partial_{\pgamma} K_{(1)}^{\idelta \jdelta}\big) K_{(\beta)}^{\pgamma \kgamma}-\big(\partial_{\pgamma} K_{(\beta)}^{\idelta \jdelta}\big) K_{(1)}^{\pgamma \kgamma}\big)=0,
\end{gather}
where $1\leq \idelta, \jdelta \leq \ldelta$, $1\leq \kgamma \leq \lgamma$ and $1\leq \delta, \gamma \leq r$. Using the expressions~(\ref{Killingtensor}) of the Killing tensors $\big(K_{(\beta)}^{\idelta \jdelta}\big)$, the fact that $\partial_{\igamma}\big(G^{\beta}\big)^{\ibeta \jbeta}=0$ for $\beta \neq \gamma$, and the fact that each of the $\lgamma \times \lgamma$ matrices $\big(\big(G^{\gamma}\big)^{\igamma \jgamma}\big)$, is invertible, we obtain that the relations~(\ref{Painlevecomm}) are equivalent to
\begin{gather}\label{Stackelcomm}
\left[\partial_{\pgamma}\left(\frac{s^{\delta 1}}{\det S}\right)\right] \frac{s^{\gamma \beta}}{\det S}-\left[\partial_{\pgamma}\left(\frac{s^{\delta \beta}}{\det S}\right)\right] \frac{s^{\gamma 1}}{\det S}=0,
\end{gather}
where $1\leq \delta, \gamma \leq r$, $1_{\gamma}\leq \pgamma \leq \lgamma $. In particular, the relations (\ref{Painlevecomm}) are independent of the block metrics (\ref{blockmetric}). Setting $\delta = \gamma$ in~(\ref{Stackelcomm}), we obtain
\[
\partial_{\pgamma}\left(\frac{s^{\gamma \beta}}{s^{\gamma 1}}\right)=0,
\]
for $1_{\gamma}\leq \pgamma \leq \lgamma$, so that using the definition (\ref{defrho}) of the quantities $\rho_{\beta \gamma}$, the relations (\ref{Stackelcomm}) take the form
\[
\partial_{\pgamma}\left(\rho_{ \beta \delta}\frac{s^{\delta 1}}{\det S}\right)\frac{s^{\gamma 1}}{\det S}=\left(\partial_{\pgamma}\left(\frac{s^{\delta 1}}{\det S}\right)\right)\rho_{\beta \gamma} \frac{s^{\gamma 1}}{\det S},
\]
which in turn reduces to (\ref{genKE}), thus proving our claim.
\end{proof}

In order to state the generalization to Painlev\'e metrics of the classical Levi-Civita conditions which hold for St\"ackel metrics, we make the hypothesis
\begin{gather}\label{H}
\rho_{\beta \delta}\neq \rho_{\beta \epsilon},\qquad \forall\, 1\leq \beta \leq r,\quad \forall\, 1\leq \delta\neq \epsilon \leq r.
\end{gather}

The generalization of the Levi-Civita conditions to the Painlev\'e case is now given by the following:
\begin{Lemma} \label{LeviCivita}
A generalized St\"ackel matrix \eqref{genStackelmat} corresponding to a Painlev\'e metric~\eqref{Painleveform} for which the genericity hypothesis~\eqref{H} holds true satisfies, for all $1\leq \beta, \gamma \leq r$,
the identities
\begin{gather}
\left(\partial_{\jalpha}\log\left(\frac{s^{\gamma 1}}{\det S}\right)\right) \left(\partial_{\kbeta}\log\left(\frac{s^{\alpha 1}}{\det S}\right)\right)+\left(\partial_{\jalpha}\log\left(\frac{s^{\beta 1}}{\det S}\right)\right)\left(\partial_{\kbeta}\log\left(\frac{s^{\gamma 1}}{\det S}\right)\right)\nonumber\\
\qquad{} -\frac{\det S}{s^{\gamma 1}}\partial_{\jalpha}\partial_{\kbeta}\left(\frac{s^{\gamma 1}}{\det S}\right)=0.\label{genLeviCivita}
\end{gather}
In particular, we have the identities
\begin{gather}\label{secondorder}
\frac{\partial}{\partial x^{\jalpha}}\frac{\partial }{\partial x^{\kbeta}}\left(\frac{\det S}{s^{\alpha1}s^{\beta1}}\right)=0,
\end{gather}
for all
$1\leq \alpha, \beta \leq r$.
\end{Lemma}
We will likewise refer to the identities (\ref{genLeviCivita}) as the {\emph{generalized Levi-Civita conditions}}.
\begin{Remark}\quad
\begin{enumerate}\itemsep=0pt
\item[1)] In the case of St\"ackel metrics, that is when $r=n$, the conditions (\ref{genLeviCivita}) reduce to the classical Levi-Civita conditions, given by
\begin{gather*}
\left(\partial_{j}\log\left(\frac{s^{l1}}{\det S}\right)\right)\left(\partial_{k}\log\left(\frac{s^{j1}}{\det S}\right)\right)+\left(\partial_{j}\log\left(\frac{s^{k1}}{\det S}\right)\right)\left(\partial_{k}\log \left(\frac{s^{l1}}{\det S}\right)\right)\\
\qquad{} -\frac{\det S}{s^{l1}}\partial_{j}\partial_{k}\left(\frac{s^{l1}}{\det S}\right)=0.
\end{gather*}
\item[2)] We shall show at the end of the next Section~\ref{PainHJSection} that the generalized Levi-Civita condi\-tions~(\ref{genLeviCivita}) hold in fact for all Painlev\'e metrics~(\ref{Painleveform}) without assuming our genericity hypothesis~(\ref{H}). Nevertheless, it is easier to obtain them from the Killing--Eisenhart equations under the assumption~(\ref{H}) as we do below.
\end{enumerate}
\end{Remark}
\begin{proof}
The general idea behind the proof is similar to the one that is used in the classical St\"ackel case, and is based on expressing the integrability conditions for the generalized Killing--Eisenhart (\ref{genKE}), with additional twist resulting from the fact that the separation is in groups of variables only. We let $1\leq \alpha \neq \beta \leq r$ denote fixed indices and introduce the simplified notation
\[
\rho_{\delta}:=\rho_{ \beta \delta}=\frac{s^{\delta \beta}}{s^{\delta 1}},
\]
so as to make the expressions a bit more compact. The generalized Killing--Eisenhart equa\-tions~(\ref{genKE}) thus take the form
\begin{gather*}%\label{simpgenKE}
\partial_{\pgamma} \rho_{\delta}= (\rho_{\gamma}-\rho_{\delta })\left(\partial_{\pgamma}\left(\log \frac{s^{\delta 1}}{\det S}\right)\right),
\end{gather*}
where $1\leq \delta,\gamma \leq r$. The integrability conditions
\begin{gather}\label{intKE}
\partial_{\kepsilon}\big(\partial_{\pgamma} \rhodelta \big)=\partial_{\pgamma}\big(\partial_{\kepsilon} \rhodelta \big)
\end{gather}
are now easily obtained. We have
\begin{gather*}
\partial_{\kepsilon} (\partial_{\pgamma} \rhodelta )= \left[(\rhoepsilon-\rhogamma)\partial_{\kepsilon}\log\left(\frac{s^{\gamma 1}}{\det S}\right)-(\rhoepsilon-\rhodelta)\partial_{\kepsilon}\log\left(\frac{s^{\delta 1}}{\det S}\right)\right] \partial_{\pgamma} \log \frac{s^{\delta 1}}{\det S} \\
\hphantom{\partial_{\kepsilon} (\partial_{\pgamma} \rhodelta )=}{} + (\rhogamma-\rhodelta)\partial_{\kepsilon}\partial_{\pgamma}\log\left(\frac{s^{\delta 1}}{\det S}\right),
\end{gather*}
so that (\ref{intKE}) becomes
\begin{gather}
(\rho_{\gamma}-\rho_{\epsilon})\left[\left(\partial_{\kepsilon}\log\left(\frac{s^{\gamma 1}}{\det S}\right)\right)\left(\partial_{\pgamma}\log\left(\frac{s^{\delta 1}}{\det S}\right)\right)\right.\nonumber\\
\qquad{} +\left(\partial_{\kepsilon}\log\left(\frac{s^{\delta 1}}{\det S}\right)\right)\left(\partial_{\pgamma}\log\left(\frac{s^{\epsilon 1}}{\det S}\right)\right) \nonumber\\
\left.\qquad{} -\left(\partial_{\kepsilon}\log\left(\frac{s^{\delta 1}}{\det S}\right)\right)\left(\partial_{\pgamma}\log\left(\frac{s^{\delta 1}}{\det S}\right)\right)-\partial_{\kepsilon}\partial_{\pgamma}\log\left(\frac{s^{\delta 1}}{\det S}\right)\right]=0,\label{compat}
\end{gather}
where $1\leq \epsilon \neq \gamma \leq r$ and $1\leq \delta \leq r$. Using the rank hypothesis (\ref{H}), expanding the logarithmic second derivative in (\ref{compat}) using the identity
\begin{gather}\label{identity}
\partial_{xy}\log f = \frac{1}{f}\partial_{x}\partial_{y}f-(\partial_{x}\log f)(\partial_{y}\log f)
\end{gather}
and
 relabeling the indices, we obtain
\begin{gather*}
\left(\partial_{\jalpha}\log\left(\frac{s^{\gamma 1}}{\det S}\right)\right) \left(\partial_{\kbeta}\log\left(\frac{s^{\alpha 1}}{\det S}\right)\right)+\left(\partial_{\jalpha}\log\left(\frac{s^{\beta 1}}{\det S}\right)\right)\left(\partial_{\kbeta}\log\left(\frac{s^{\gamma 1}}{\det S}\right)\right)\\
\qquad{}-\frac{\det S}{s^{\gamma 1}}\partial_{\jalpha}\partial_{\kbeta}\left(\frac{s^{\gamma 1}}{\det S}\right)=0,
\end{gather*}
which is precisely (\ref{genLeviCivita}). Finally, the relations~(\ref{secondorder}) are obtained by setting $\delta = \epsilon$ in the integrability condition (\ref{compat}), using the identity (\ref{identity}), and the fact that the cofactors $s^{\gamma \alpha}$ and $s^{\epsilon \alpha}$ are independent of the groups of variables ${\bf x}^{\gamma}$ and ${\bf x}^{\epsilon}$ respectively.
\end{proof}

\section{Different characterizations of Painlev\'e metrics}\label{PainHJSection}

Let us start with the characterization of Painlev\'e metrics in terms of complete additive separability of the Hamilton--Jacobi equations stated in Proposition \ref{HJSeparability-Painleve}.

\begin{proof}[Proof of Proposition \ref{HJSeparability-Painleve}]

In block orthogonal coordinates $(\textbf{x}^\alpha)$ associated to the metric
\[
 g = \sum_{\alpha = 1}^r c_\alpha G_\alpha,
\]
the Hamilton--Jacobi equation (\ref{HJ}) reads
\begin{gather} \label{HJ1}
 \sum_{\alpha =1}^{r} c^\alpha \sum_{\ialpha =1_{\alpha}}^{\lalpha}\sum_{\jalpha =1_{\alpha}}^{\lalpha} \big(G^{\alpha}\big)^{\ialpha \jalpha}\frac{\partial W}{\partial x^{\ialpha} }\frac{\partial W}{\partial x^{\jalpha}}=E=:a_{1},
\end{gather}
where $c^\alpha = (c_\alpha)^{-1}$. Assume that there exists a solution $W$ in the block-separable form~(\ref{Painsepform}) satisfying the completeness condition~(\ref{rankPainHJ}). Then (\ref{HJ1}) can be written as
\begin{gather} \label{HJ2}
 \sum_{\alpha =1}^{r} c^\alpha \sum_{\ialpha =1_{\alpha}}^{\lalpha}\sum_{\jalpha =1_{\alpha}}^{\lalpha} \big(G^{\alpha}\big)^{\ialpha \jalpha}\frac{\partial W_\alpha}{\partial x^{\ialpha} }\frac{\partial W_\alpha}{\partial x^{\jalpha}}=E=:a_{1}.
\end{gather}
Differentiating (\ref{HJ2}) with respect to $a_\beta$, we get
\begin{gather} \label{HJ3}
 \sum_{\alpha =1}^{r} c^\alpha \sum_{\ialpha =1_{\alpha}}^{\lalpha}\sum_{\jalpha =1_{\alpha}}^{\lalpha} 2 \big(G^{\alpha}\big)^{\ialpha \jalpha}\frac{\partial W_\alpha}{\partial x^{\ialpha} } \frac{\partial^2 W_\alpha}{\partial a_\beta \partial x^{\jalpha}} = \delta_{1\beta} .
\end{gather}
From (\ref{Painsepform}) and (\ref{rankPainHJ}), it follows that the family of matrices
\begin{gather*} %\label{SM}
 S(a_1,\dots,a_r) = (S_{\alpha \beta})(a_1,\dots,a_r) := 2 \big(G^{\alpha}\big)^{\ialpha \jalpha}\frac{\partial W_\alpha}{\partial x^{\ialpha} } \frac{\partial^2 W_\alpha}{\partial a_\beta \partial x^{\jalpha}}
\end{gather*}
are non-singular st\"ackel matrices of rank $r$ for all $(a_1,\dots,a_r) \in U$. Using the invertibility of $S(a_1,\dots,a_r)$ which is equivalent to the completeness condition (\ref{rankPainHJ}), we get immediately from~(\ref{HJ3}) that
\[
 c^\alpha = \frac{s^{\alpha 1}}{\det S},
\]
which proves that $g$ is a Painlev\'e metric.

Conversely, assume that $g$ is a Painlev\'e metric of the form (\ref{Painleveform}). Then the Hamilton--Jacobi equation (\ref{HJ}) takes the form
\begin{gather}\label{HJPainform}
\sum_{\alpha =1}^{r} \left(\frac{s^{\alpha 1}}{\det S}\right) \sum_{\ialpha =1_{\alpha}}^{\lalpha}\sum_{\jalpha =1_{\alpha}}^{\lalpha} \big(G^{\alpha}\big)^{\ialpha \jalpha}\frac{\partial W}{\partial x^{\ialpha} }\frac{\partial W}{\partial x^{\jalpha}}=E=:a_{1}.
\end{gather}
Now, since
\[
\sum_{\alpha =1}^{r}\left(\frac{s^{\alpha 1}}{\det S}\right)s_{\alpha \beta}=\delta_{1 \beta},
\]
we may rewrite (\ref{HJPainform}) as
\begin{gather} \label{HJ4}
\sum_{\alpha =1}^{r}\frac{s^{\alpha 1}}{\det S}\left(\sum_{\ialpha =1_{\alpha}}^{\lalpha}\sum_{\jalpha =1_{\alpha}}^{\lalpha}\big(G^{\alpha}\big)^{\ialpha \jalpha}\frac{\partial W}{\partial x^{\ialpha} }\frac{\partial W}{\partial x^{\jalpha}}-\sum_{\beta=1}^{r}s_{\alpha \beta}a_{\beta}\right)=0,
\end{gather}
for any $(a_1,\dots,a_r) \in U \subset \mathbb{R}^r$. Choosing $W$ in the block-separable form~(\ref{Painsepform}), we see that any solution of the reduced Hamilton--Jacobi equations
\begin{gather} \label{HJ5}
\sum_{\ialpha =1_{\alpha}}^{\lalpha}\sum_{\jalpha =1_{\alpha}}^{\lalpha}\big(G^{\alpha}\big)^{\ialpha \jalpha}\frac{\partial W_{\alpha}}{\partial x^{\ialpha} }\frac{\partial W_{\alpha}}{\partial x^{\jalpha}}=\sum_{\beta = 1}^{r}s_{\alpha \beta}a_{\beta},
\end{gather}
will provide a solution of (\ref{HJ4}). But the latter equation always admit \emph{locally} solutions $W_\alpha(\textbf{x}^\alpha, \allowbreak a_1,\dots,a_r)$ by standard PDE results \cite{Tay2011a}. Differentiating (\ref{HJ5}) with respect to $a_{\beta}$, we obtain
\begin{gather*} %\label{HJ6}
2\sum_{\alpha =1}^{r}\sum_{\ialpha =1_{\alpha}}^{\lalpha}\sum_{\jalpha =1_{\alpha}}^{\lalpha}\big(G^{\alpha}\big)^{\ialpha \jalpha}\left(\frac{\partial W_{\alpha}}{\partial x^{\ialpha}}\right)\left( \frac{\partial^{2}W_{\alpha}}{\partial a_{\beta}\partial x^{\jalpha}}\right)=s_{\alpha \gamma}.
\end{gather*}
Since the generalized St\"ackel matrix (\ref{genStackelmat}) is to be non-singular, it follows that the rank condition that must be satisfied by our block-separable parametrized family of solutions of the Hamilton--Jacobi equation is precisely~(\ref{rankPainHJ}), thus proving that metrics of the Painlev\'e form~(\ref{Painleveform}) are indeed characterized by the existence of a parametrized family of solutions of the Hamilton--Jacobi equation satisfying a suitable completeness condition.
\end{proof}

Let us now give another proof of Proposition~\ref{HJSeparability-Painleve} that will allow us to characterize Painlev\'e metrics in terms of the generalized Levi-Civita conditions (\ref{genLeviCivita}). Working in block-orthogonal coordinates or directly with a Painlev\'e metric with the identification
\begin{gather}\label{not}
c^\alpha:=\frac{\det S}{s^{\alpha 1}},
\end{gather}
the Hamilton--Jacobi equation (\ref{HJ}) takes the form
\begin{gather}\label{HJPainformnew}
\sum_{\alpha =1}^{r} c^\alpha \sum_{\ialpha =1_{\alpha}}^{\lalpha}\sum_{\jalpha =1_{\alpha}}^{\lalpha} \big(G^{\alpha}\big)^{\ialpha \jalpha}\frac{\partial W}{\partial x^{\ialpha} }\frac{\partial W}{\partial x^{\jalpha}}=E=:a_{1}.
\end{gather}
We recall that we seek a solution $W$ of (\ref{HJPainformnew}) which is additively separable into groups of variables, that is
\[
W=\sum_{\alpha =1}^{r}W_{\alpha}\big({\bf x}^{\alpha}\big),
\]
and let
\[
u_{\alpha}^{(1)}=\sum_{\ialpha =1_{\alpha}}^{\lalpha}\sum_{\jalpha =1_{\alpha}}^{\lalpha}\big(G^{\alpha}\big)^{\ialpha \jalpha}\frac{\partial W_{\alpha}}{\partial x^{\ialpha} }\frac{\partial W_{\alpha}}{\partial x^{\jalpha}},
\]
in which case the Hamilton--Jacobi equation (\ref{HJPainformnew}) takes the form
\[
\sum_{\alpha =1}^{r}c^{\alpha}u_{\alpha}^{(1)}=a_{1}.
\]
Differentiating the latter equation with respect to $x^{\kbeta}$, we obtain the first order differential system in normal form
\begin{gather}\label{normalsyst}
 \partial_{\kbeta}u^{(1)}_{\gamma} =   \begin{cases}\displaystyle -\frac{1}{c^{\beta}}\sum_{\alpha =1}^{r}\big(\partial_{\kbeta} c^{\alpha}\big) u^{(1)}_{\alpha} & \gamma = \beta, \\
 0 & \gamma \neq \beta . \end{cases}
\end{gather}
Introducing the first-order differential operators
\begin{gather}\label{totaldiff}
D_{\kbeta} := \frac{\partial}{\partial x^{\kbeta}}-\frac{1}{c^{\beta}}\sum_{\alpha =1}^{r}\big(\partial_{\kbeta} c^{\alpha}\big) u^{(1)}_{\alpha}\frac{\partial}{\partial u^{(1)}_{\beta}},
\end{gather}
the differential system (\ref{normalsyst}) will admit a family of solutions $u^{(1)}_{\alpha}=u^{(1)}_{\alpha}\big(x^{1},\dots,x^{n};a_{1},\dots,a_{r}\big)$ defined on an open subset $U\in M$ and depending smoothly on~$r$ constants $(a_{1},\dots,a_{r})$ defined in an open subset $A\subset {\mathbb{R}}^{r}$, satisfying the rank condition
\begin{gather}\label{rankHJ}
\det\left(\frac{\partial u^{(1)}_{\alpha}}{\partial a_{\beta}}\right)\neq 0,
\end{gather}
if and only if the operators (\ref{totaldiff}) pairwise commute, that is
\begin{gather}\label{commutation}
\big[D_{\kbeta},D_{\jalpha}\big]=0
\end{gather}
for all $1\leq \alpha, \beta \leq r$, $1_{\alpha}\leq \jalpha\leq \lalpha$, $1_{\beta}\leq \kbeta\leq \lbeta$. We refer to \cite[Theorem~2.1]{BCR1-2002} for this result. Note that if~(\ref{commutation}) hold, then the Hamilton--Jacobi equation admits locally a solution which is additively separable into groups of variables and satisfies the completeness condition~(\ref{rankPainHJ}) as a consequence of the completeness condition (\ref{rankHJ}). In consequence, such metrics are of the Painlev\'e form~(\ref{Painleveform}).

We now prove that the commutation relations (\ref{commutation}) are equivalent to the generalized Levi-Civita conditions~(\ref{genLeviCivita}). Note that this is a natural generalization of the link between the complete separation of variables which is familiar from the St\"ackel case and the classical Levi-Civita separability conditions, as reviewed in~\cite{BCR1-2002,KaMi1984}. Indeed, we have
\begin{Lemma}
The pairwise commutation relations \eqref{commutation} for the derivations $D_{\kbeta}$ are equivalent to the generalized Levi-Civita conditions~\eqref{genLeviCivita}.
\end{Lemma}

\begin{proof}Substituting the expression (\ref{totaldiff}) of the differential operators $D_{\kbeta}$ into the commutation relations (\ref{commutation}), we obtain
 \begin{gather*}
 \left[-\partial_{\kbeta}\left(\frac{1}{c^{\gamma}}\sum_{\delta =1}^{r}\big(\partial_{\pgamma} c^{\delta}\big)u^{(1)}_{\delta}\right) + \frac{1}{c^{\beta} c^{\gamma}}\sum_{\alpha =1}^{r}\big(\partial_{\kbeta} c^{\alpha}\big)u^{(1)}_{\alpha}\big(\partial_{\pgamma} c^{\beta}\big)\right]\frac{\partial}{\partial u^{(1)}_{\gamma}} \\
 \qquad{} +
 \left[-\partial_{\kbeta}\left(\frac{1}{c^{\gamma}}\sum_{\delta =1}^{r}\big(\partial_{\pgamma} c^{\delta}\big)u^{(1)}_{\delta}\right)+\frac{1}{c^{\beta}c^{\gamma}}\sum_{\alpha =1}^{r}\big(\partial_{\kbeta}c^{\alpha}\big)u^{(1)}_{\alpha}\big(\partial_{\pgamma} c^{\beta}\big)\right]\frac{\partial}{\partial u^{(1)}_{\gamma}}=0.
 \end{gather*}
For $1\leq \gamma \neq \beta \leq r$, the above identity is equivalent to
\begin{gather} \label{GenLC}
\big(\partial_{\kbeta} \log c^{\gamma}\big)\big(\partial_{\pgamma} \log c^{\alpha}\big)+\big(\partial_{\kbeta} \log c^{\alpha}\big)\big(\partial_{\pgamma} \log c^{\beta}\big)-\frac{1}{c^{\alpha}}\partial_{\kbeta}\partial_{\pgamma}c^{\alpha}=0,
\end{gather}
which are precisely the generalized Levi-Civita conditions~(\ref{genLeviCivita}) after relabeling of indices. Note that for $1\leq \gamma = \beta \leq r$, the above identity is always satisfied by a straightforward calculation.
\end{proof}

At this stage, we thus have proved another characterization of Painlev\'e metrics which appears in \cite[p.~12]{CR2019}.

\begin{Proposition} \label{PainLC}A metric $g$ is of Painlev\'e type if and only if there exist block orthogonal coordinates~$\big(\textbf{x}^\alpha\big)$ such that the generalized Levi-Civita conditions \eqref{GenLC} hold.
\end{Proposition}

We finish this section giving still another characterization of Painlev\'e metrics of a more intrinsic nature. The starting point is the observation that the generalized Killing--Eisenhart equations are related to the existence of quadratic first integrals (or symmetries) $K_{(\beta)}$ by the following result proved in~\cite[Proposition~5.3]{CR2019} (see also Lemma~\ref{KE} for an implicit proof of this proposition)

\begin{Proposition}In block orthogonal coordinates, we have that
\[
 \big\{H,K_{(\beta)} \big\} = 0,
\]	
if and only if the Killing--Eisenhart equations
\begin{gather}\label{genKE1}
 \partial_{j_{\gamma}}(\rho_{ \beta \delta})= (\rho_{\beta \gamma} - \rho_{\beta \delta }) \big(\partial_{j_{\gamma}} \log c^\delta\big),
\end{gather}
hold for all $1 \leq \gamma, \delta \leq r$.
\end{Proposition}

The second observation is the fact that the integrability conditions for the Killing--Eisenhart equations~(\ref{genKE1}) are given by
\begin{gather*} %\label{IC1}
 (\rho_{\beta \gamma} - \rho_{\beta \delta}) \left[ \big( \partial_{\jgamma} \log c^{\delta} \big) \big( \partial_{\kalpha} \log c^{\gamma}\big) + \big(\partial_{\jgamma} \log c^{\alpha} \big) \big( \partial_{\kalpha} \log c^{\delta} \big) - \frac{1}{c^{\delta}} \partial_{\kalpha}\partial_{\jgamma}c^{\delta} \right] = 0,
\end{gather*}
for all $1 \leq \alpha, \beta, \gamma, \delta \leq r$. Clearly, these integrability conditions can be shortened as
\begin{gather} \label{IC2}
 (\rho_{\beta \gamma} - \rho_{\beta \delta}) \cdot \left[ \textrm{Levi-Civita conditions} \right] = 0.
\end{gather}

Using these two observations, Chanu ad Rastelli proved in \cite[Proposition 5.5]{CR2019} the following characterization of Painlev\'e metrics.

\begin{Proposition} \label{PainKA}
$(M,g)$ is a Painlev\'e manifold if and only if
\begin{enumerate}\itemsep=0pt
\item[$1)$] there exist $r$ independent quadratic first integral $K_{(\beta)}$, $\beta = 1,\dots,r$ such that $K_{(1)} = H$ and $\{ H, K_{(\beta)} \} = 0$.
\item[$2)$] The associated Killing two tensors	$K_{(\beta)}^{ij}$ are simultaneously block-diagonalized and have common normally integrable eigenspaces.
\end{enumerate}
\end{Proposition}
\begin{proof}
If $g$ is a Painlev\'e metric, then the above assertions were already proved.

Assume now that there exist $r$ linearly independent Killing tensors simultaneously in block diagonal form. Then the generalized Killing--Eisenhart equations (\ref{genKE1}) will admit an $r$-dimen\-sio\-nal vector space of solutions. The latter is equivalent to the invertibility of the fundamental matrix $\mathcal{A}$ defined by
\[
{\mathcal{A}}=\left(
\begin{matrix}
\rho_{11}&\dots&\rho_{1r} \\
\vdots& & \vdots \\
\rho_{r1}&\dots&\rho_{rr}
\end{matrix}
\right).
\]
Hence any solution $(\rho_{1},\dots,\rho_{r})$ of (\ref{genKE1}) is given by
\[
\left(
\begin{matrix}\rho_{1}\\
\vdots\\
\rho_{r}
\end{matrix}
\right)
=
{\mathcal{A}}\left(
\begin{matrix}a_{1}\\
\vdots\\
a_{r}
\end{matrix}
\right),
\]
for some constants $(a_{1},\dots,a_{r})\in A$. It is clear then that the Killing--Eisenhart equations are completely integrable and thus satisfy the integrability conditions (\ref{IC2}). Moreover, from the invertibility of~$\mathcal{A}$, we can always choose the constants $(a_1,\dots,a_r)$ such that $\rho_{\alpha} \neq \rho_{\beta}$, $\forall\, 1 \leq \alpha \ne \beta \leq r$ at a point $p\in M$ and therefore in an neighbourhood of $p$, by continuity. Hence, the integrabi\-li\-ty conditions~(\ref{IC2}) reduces to the generalized Levi-Civita separability conditions~(\ref{GenLC}). We conclude that $g$ is a Painlev\'e metric from Proposition~\ref{PainLC}.
\end{proof}

As a concluding remark for this section, we emphasize that the hypotheses~(\ref{H}) that we make to deduce the Levi-Civita conditions from the Killing--Eisenhart equations aren't in fact necessary. Indeed, it follows from Proposition~\ref{PainLC} or Proposition~\ref{PainKA} that the Levi-Civita conditions always hold whenever~$g$ is a Painlev\'e metric, that is a metric of the form~(\ref{Painleveform}).

